%--------------------------------------------------------------------------------
%
%
%--------------------------------------------------------------------------------
\pagebreak
\section*{MBR move and branch}
\addcontentsline{toc}{section}{MBR move and branch}

\instructionentry{Syntax}{

	 \texttt{MBR.<cond> RegR, RegB, target}
    
   	where <cond> is EQ, LT, NE or LE.
}

\instructionentry{Format}{

	The MBR instruction uses the xxx  instruction format. \\

   	\begin{flushleft}
	\begin{tikzpicture}[scale=0.7, transform shape]
        
        % \draw[help lines, gray!70, dashed] (0,0) grid(16,1.5);
        
        \node[  tsRectangle, 
                minimum width=16cm,
                minimum height=1cm,
                text centered,
                fill=white!50] (blocks) at (8,1) { };
                
     	\draw[-] (1,0.5) -- (1, 1.5);
     	\draw[-] (3,0.5) -- (3, 1.5);
     	\draw[-] (5,0.5) -- (5, 1.5);
     	\draw[-] (6.5,0.5) -- (6.5, 1.5);
     	\draw[-] (8.5,0.5) -- (8.5, 1.5);
          	
     	\node at (0.5,0.25) {2};
    	\node at (2,0.25) {4};
    	\node at (4,0.25) {4};
    	\node at (5.75,0.25) {3};
    	\node at (7.5,0.25) {4};
   	 	\node at (11.5,0.25) {15};
   	 	
   	 	\node at (0.5,1) {\small 0};
		\node at (2,1) {\small \OpCodeMBR};
		\node at (4,1) {\small Reg R};
		\node at (5.75,1) {\small 0};
		\node at (7.5,1) {\small Reg B};
		\node at (11.5,1) {\small ofs};
   	 	
       \end{tikzpicture}	\end{flushleft}
}

\instructionentry{Description}{

The \texttt{MBR} instruction ...}

\instructionentry{Operation}{

  	\texttt{RegR <- RegR + ofs; }
} 

\instructionentry{Exceptions}{

	None.
}

\instructionentry{Notes}{

	None. 
}




